\subsection{Primer cuello de botella bajo la carga de septiembre}
\label{sec:cuello-botella-septiembre}

Septiembre eleva las entregas diarias de aproximadamente 1.400 a 2.600 durante tres semanas. Las 2.600 entregas ocurren a lo largo de la jornada; la ventana 05:30--07:00 concentra la salida de 96 camiones, no todas las firmas de entrega. En esa ventana, el primer cuello de botella probable es la confirmación de guías del ERP mediante la ACL: es una dependencia singular y aumentar réplicas Laravel no acelera al emisor legado. La base transaccional de Talca, la WAN al reconectar y la transferencia de evidencias son candidatos adicionales, no recursos cuya saturación ya se haya medido. La Tabla~\ref{tab:cuello-botella} muestra cómo contrastar la hipótesis; RT-09.05 exige sustentar el primer límite con una prueba reproducible.

\begin{tablalafrox}{Detección del primer cuello de botella}{tab:cuello-botella}{F{0.23} F{0.26} F{0.27} Y}{\cab{Recurso} & \cab{Señal de saturación} & \cab{Respuesta lógica} & \cab{Prueba}}
ERP por ACL & Aumentan guías sin confirmar y latencia p95 de emisión. & Preemitir al cerrar la carga nocturna; invalidar y reemitir si cambia. & 96 camiones en ventana. \\
VM-02, base transaccional de Talca & Crecen la latencia de confirmación y la espera de escritura. & Medir y aliviar contención de escrituras del WMS local. & Carga de preparación y despacho de Talca. \\
Cola de integración & Crecen edad del mensaje y DLQ. & Priorizar guías; diferir reportes. & Corte y reconexión. \\
M2 en bodega & Crecen espera de reserva y bloqueos de stock. & Serializar por SKU y aislar consultas. & Doble compromiso en peak. \\
\end{tablalafrox}

{\footnotesize Fuente: elaboración propia de LafroX a partir de Caso 02, numeral 14.1 y ventana 05:30--07:00.\par}

La emisión de guías debe medirse antes de atribuir el límite al cómputo escalable.

Al cerrar la carga nocturna, \texttt{erp-sync} en VM-04 solicita la guía al ERP mediante la ACL local; el ERP dispone de fibra, LTE y satélite para comunicarse con el SII. Cualquier cambio de carga invalida la guía y exige otra antes de liberar. Se mide el número de guías sin confirmar antes de cada salida y se alerta al responsable de despacho; el camión afectado no se marca como liberado mientras falte el documento o una contingencia autorizada. La prueba de carga debe confirmar si el ERP es efectivamente el primer límite y ajustar esta hipótesis con medición, como exige RT-09.05.

El Anexo 4-T especifica percentiles, ventanas de ensayo, carga de 1,5 veces el peak y escenario de crecimiento de tres veces, con colas y procesos aislados. Los resultados se deben medir; los cálculos de escenario no son evidencia de capacidad ejecutada.


